\section{Прямая сумма подпространств. Теорема о разложении ЛП в прямую сумму подпространств}

\begin{definition}
Сумма подпространств $U_1 + U_2$ называется \textbf{прямой} (обозначается $U_1 \oplus U_2$), если любой вектор $x \in U_1 + U_2$ единственным образом представляется в виде суммы $x = x_1 + x_2$, где $x_1 \in U_1$ и $x_2 \in U_2$.
\end{definition}

\begin{theorem}[Критерий прямой суммы]
Сумма $U_1 + U_2$ является прямой тогда и только тогда, когда пересечение этих подпространств состоит только из нулевого вектора:
\[
U_1 \cap U_2 = \{0\}.
\]
\end{theorem}

\begin{proof}
\begin{itemize}
    \item \textbf{Необходимость ($\Rightarrow$):} Пусть сумма прямая. Предположим, что существует ненулевой вектор $v \in U_1 \cap U_2$. Тогда нулевой вектор можно представить двумя разными способами:
    \[
    0 = 0 + 0 \quad \text{и} \quad 0 = v + (-v),
    \]
    где $v \in U_1, -v \in U_2$. Это противоречит единственности разложения. Следовательно, $U_1 \cap U_2 = \{0\}$.
    \item \textbf{Достаточность ($\Leftarrow$):} Пусть $U_1 \cap U_2 = \{0\}$. Допустим, вектор $x$ имеет два разложения: $x = x_1 + x_2 = x_1' + x_2'$. Тогда:
    \[
    x_1 - x_1' = x_2' - x_2.
    \]
    Левая часть принадлежит $U_1$, правая --- $U_2$. Значит, этот вектор принадлежит пересечению $U_1 \cap U_2$. Но пересечение содержит только ноль, поэтому $x_1 - x_1' = 0$ и $x_2' - x_2 = 0$, то есть $x_1 = x_1'$ и $x_2 = x_2'$. Разложение единственно.
\end{itemize}
\hfill$\blacksquare$
\end{proof}

\begin{theorem}[О разложении пространства в прямую сумму]
Пусть $V$ --- конечномерное линейное пространство, и $V = U_1 \oplus U_2 \oplus \dots \oplus U_k$ (прямая сумма подпространств). Если в каждом подпространстве $U_i$ выбрать по базису, то объединение этих базисов образует базис всего пространства $V$.
\end{theorem}

\begin{proof}
Пусть $E_i = \{e_{i1}, \dots, e_{im_i}\}$ --- базис подпространства $U_i$. Рассмотрим объединенную систему $E = E_1 \cup E_2 \cup \dots \cup E_k$.
\begin{enumerate}
    \item \textit{Система $E$ порождает $V$:} Любой вектор $x \in V$ представим как $x = x_1 + \dots + x_k$, где $x_i \in U_i$. Каждый $x_i$ линейно выражается через базис $E_i$. Следовательно, $x$ линейно выражается через объединение $E$.
    \item \textit{Система $E$ линейно независима:} Составим линейную комбинацию всех векторов системы $E$, равную нулю:
    \[
    \sum_{i=1}^k \left( \sum_{j=1}^{m_i} \lambda_{ij} e_{ij} \right) = 0.
    \]
    Обозначим $v_i = \sum_{j=1}^{m_i} \lambda_{ij} e_{ij}$. Очевидно, что $v_i \in U_i$. Тогда мы получили разложение нуля: $v_1 + v_2 + \dots + v_k = 0$. 
    Так как сумма $U_1 \oplus \dots \oplus U_k$ прямая, разложение любого вектора (в том числе нуля) единственно. Единственное разложение нуля --- это $0 + 0 + \dots + 0$. Следовательно, $v_i = 0$ для всех $i = 1, \dots, k$.
    Но $v_i$ --- это линейная комбинация базисных векторов подпространства $U_i$. В силу линейной независимости базиса $E_i$, все коэффициенты $\lambda_{ij} = 0$. 
\end{enumerate}
Таким образом, $E$ --- линейно независимая система, порождающая $V$, то есть базис $V$.
\hfill$\blacksquare$
\end{proof}

\begin{corollary}
Если $V = U_1 \oplus U_2$, то $\dim V = \dim U_1 + \dim U_2$. Это частный случай формулы Грассмана, так как $\dim(U_1 \cap U_2) = \dim\{0\} = 0$.
\end{corollary}
